%!TEX program = xelatex

%%%%%%%%%%%%%%%%%%%%%%%%%%%%%%%%%%%%%%%%%
% Thin Sectioned Essay
% LaTeX Template
% Version 1.0 (3/8/13)
%
% This template has been downloaded from:
% http://www.LaTeXTemplates.com
%
% Original Author:
% Nicolas Diaz (nsdiaz@uc.cl) with extensive modifications by:
% Vel (vel@latextemplates.com)
%
% License:
% CC BY-NC-SA 3.0 (http://creativecommons.org/licenses/by-nc-sa/3.0/)
%
%%%%%%%%%%%%%%%%%%%%%%%%%%%%%%%%%%%%%%%%%

%----------------------------------------------------------------------------------------
%	PACKAGES AND OTHER DOCUMENT CONFIGURATIONS
%----------------------------------------------------------------------------------------

\documentclass[a4paper, 11pt]{article} % Font size (can be 10pt, 11pt or 12pt) and paper size (remove a4paper for US letter paper)
\usepackage[table]{xcolor}
\usepackage{media9}
\usepackage{algorithm,algpseudocode}
\usepackage{algorithmicx}
\definecolor{keycolor}{RGB}{172, 42, 42}
\definecolor{mbleu}{RGB}{64,96,127}
\definecolor{vimvert}{RGB}{46, 139, 87}
\definecolor{mygreen}{rgb}{0,0.6,0}
\definecolor{mygray}{rgb}{0.5,0.5,0.5}
\definecolor{mymauve}{rgb}{0.58,0,0.82}
\usepackage{xeCJK}
% \setCJKmainfont[AutoFakeSlant=0.25]{SourceHanSerifSC-Regular}
% \setCJKsansfont[AutoFakeSlant=0.25]{SourceHanSansSC-Regular}
% \setCJKmonofont[AutoFakeSlant=0.25]{Sarasa Mono SC}

\usepackage{geometry}
\geometry{left=2.54cm, top=2.54cm, right=2.54cm, bottom=2.54cm}
\usepackage{subcaption}
\usepackage{tikz}
\usetikzlibrary{tikzmark}
\usepackage{listings}
\usepackage{color}
\usepackage{forest}
\usepackage{float}
\usepackage{makecell}
\usepackage{url}
\usepackage{siunitx}
\usepackage{amsmath}
\usepackage{listings} 
\lstset{language=C++}%这条命令可以让LaTeX排版时将C++键字突出显示
\lstset{breaklines}%这条命令可以让LaTeX自动将长的代码行换行排版
\lstset{extendedchars=ture}
\lstset{
    columns=fixed,       
    numbers=left,                                        % 在左侧显示行号
    frame=none,                                          % 不显示背景边框
    backgroundcolor=\color[RGB]{245,245,244},            % 设定背景颜色
    keywordstyle=\color[RGB]{40,40,255},                 % 设定关键字颜色
    numberstyle=\footnotesize\color{darkgray},           % 设定行号格式
    commentstyle=\it\color[RGB]{0,96,96},                % 设置代码注释的格式
    stringstyle=\rmfamily\slshape\color[RGB]{128,0,0},   % 设置字符串格式
    showstringspaces=false,                              % 不显示字符串中的空格
    language=C++,                                        % 设置语言
}
% \usepackage{fancyhdr}
% \usepackage{lastpage}
% \pagestyle{fancy}
% \fancyhf{}
% \fancyfoot[L]{}
% \fancyfoot[R]{第 \thepage 页，共 \pageref{LastPage} 页}
% % \fancyfoot[C]{\thepage/\pageref{LastPage}}
% \renewcommand{\headrulewidth}{0pt} 
% \renewcommand{\footrulewidth}{0.4pt} 

% \usepackage[protrusion=true,expansion=true]{microtype} % Better typography
\usepackage{graphicx} % Required for including pictures
\usepackage{wrapfig} % Allows in-line images
\usepackage{newfloat}
\usepackage{amsmath}
\usepackage{hyperref}
\usepackage{listings}
\usepackage{geometry}
\usepackage{setspace}
\usepackage{mathpazo} % Use the Palatino font
\usepackage[T1]{fontenc} % Required for accented characters
\linespread{1.2} % Change line spacing here, Palatino benefits from a slight increase by default
% \setlength{\parskip}{0.2em}
\usepackage{amsmath}
\newtheorem{theorem}{{定理}}
\newtheorem{lemma}{{推论}}
\newtheorem{proof}{Proof}[section]
\newtheorem{Def}{{定义}}
\newtheorem{proposition}{{命题}}
\usepackage{indentfirst}

\usepackage{pdfpages}
\setlength{\parindent}{2em}

\makeatletter
\renewcommand\@biblabel[1]{\textbf{#1.}} % Change the square brackets for each bibliography item from '[1]' to '1.'
\renewcommand{\@listI}{\itemsep=0pt} % Reduce the space between items in the itemize and enumerate environments and the bibliography

\renewcommand{\maketitle}{ % Customize the title - do not edit title and author name here, see the TITLE block below
\begin{center} % Right align
{\LARGE\@title} % Increase the font size of the title

\large{\@subtitle}

\vspace{1em} % Some vertical space between the title and author name

{\large\@author} % Author name
% \\\@date % Date

% \vspace{1.5em} % Some vertical space between the author block and abstract
\end{center}
}

\renewcommand{\figurename}{图}
\renewcommand{\tablename}{表}

%----------------------------------------------------------------------------------------
%	TITLE
%----------------------------------------------------------------------------------------
\usepackage{geometry}
\geometry{verbose,letterpaper}
\usepackage{movie15}
\usepackage{hyperref}
\usepackage{enumitem,amssymb}
\newlist{todolist}{itemize}{2}
\setlist[todolist]{label=$\square$}
% for check symbol 
\usepackage{pifont}
\newcommand{\cmark}{\ding{51}}%
\newcommand{\xmark}{\ding{55}}%
\newcommand{\done}{\rlap{$\square$}{\raisebox{2pt}{\large\hspace{1pt}\cmark}}\hspace{-2.5pt}}
\newcommand{\wontfix}{\rlap{$\square$}{\large\hspace{1pt}\xmark}}

\begin{document}
    
\includegraphics[scale=0.3]{sysu.jpeg}

\title{\textbf{\huge Leetcode}\\ % Title
\textbf{  }\\
\textbf{\huge 题解}
} % Subtitle
\newcommand\@subtitle{}%Matrix 6382题}


    
    
 % Institution



%----------------------------------------------------------------------------------------


\maketitle % Print the title section


\begin{center}

姓名：刘斯宇\\
学号：17341110
\end{center}

\tableofcontents
%----------------------------------------------------------------------------------------
%	ABSTRACT AND KEYWORDS
%----------------------------------------------------------------------------------------
\newpage


% \setcounter{tocdepth}{3}
% \renewcommand{\contentsname}{目录}
% \tableofcontents

% \vspace{2em} % Some vertical space between the abstract and first section

\newpage


\begin{spacing}{1.5}

\section{Two Sum}
\section*{\href{https://leetcode.com/problems/two-sum/}{题目}}
\fbox{%
  \parbox{\textwidth}{%

  \textbf{Given an array of integers, return indices of the two numbers such that they add up to a specific target.}

  \textbf{You may assume that each input would have exactly one solution, and you may not use the same element twice.}
  }
  }
  \section*{例子}
  \fbox{%
  \parbox{\textwidth}{%
  Given nums = [2, 7, 11, 15], target = 9,

Because nums[0] + nums[1] = 2 + 7 = 9,
return [0, 1].
  }
  }

\section*{思路}
\textbf{\color{blue}配对的题目，我们当然可以用两层额for循环来解决这个问题，但是这样的话，就没什么意义了，所以我们要降低时间复杂度，也就是要降低内层查找的时间复杂度，如何降低？用空间来换时间，所以很显然需要有个map，这样就很容易的得到思路了。代码中比较巧妙的一点是如果出现了vector有相同的值的话，但是我们用后面的induce去覆盖前面的，这样我们遍历的时候就不会忘掉这种相同的值，但是induce不同的情况。}
\begin{lstlisting}[caption={},frame=shadowbox]
vector<int> twoSum(vector<int>& nums, int target)
{
    unordered_map<int,int>value2index;
    for(int i=0;i<nums.size();++i)
        value2index[nums[i]]=i;
    for(int i=0;i<nums.size();++i)
    {
        int temp=target-nums[i];
        if(value2index.find(temp)!=value2index.end()&&i!=value2index[temp])
            return {i,value2index[temp]};
    }
    return {};
}   
\end{lstlisting}

\section{Add Two Numbers}
\section*{\href{https://leetcode.com/problems/add-two-numbers/}{题目} }
\fbox{%
  \parbox{\textwidth}{%

  \textbf{You are given two non-empty linked lists representing two non-negative integers. The digits are stored in reverse order and each of their nodes contain a single digit. Add the two numbers and return it as a linked list.}

  \textbf{You may assume the two numbers do not contain any leading zero, except the number 0 itself.}
  }
  }
  \section*{例子}
  \fbox{%
  \parbox{\textwidth}{%
  Input: (2 -> 4 -> 3) + (5 -> 6 -> 4)\\
  Output: 7 -> 0 -> 8\\
  Explanation: 342 + 465 = 807.
  }
  }

\section*{思路}
\textbf{\color{blue}链表的题目比较简单，一直向前遍历，直到找到结束的标志--两个链表都遍历完了以及没有进位了。}
\begin{lstlisting}[caption={},frame=shadowbox]
    ListNode* addTwoNumbers(ListNode* l1, ListNode* l2)
    {
        ListNode*temp1=l1;
        ListNode*temp2=l2;
        ListNode*start=new ListNode(0);
        ListNode*mm=start;
        int jinwei=0;
        while(temp1||temp2||jinwei)
        {
            int sum=0;
            if(temp1)
            {
                sum+=temp1->val;
                temp1=temp1->next;
            }
            if(temp2)
            {
                sum+=temp2->val;
                temp2=temp2->next;
            }
            if(jinwei)
                sum+=1;
            if(sum>=10)
                jinwei=1;
            else
                jinwei=0;
            sum=sum%10;
            ListNode*temp=new ListNode(sum);
            start->next=temp;
            start=temp;
        }
        return mm->next;
    }
\end{lstlisting}

\section{ Longest Substring Without Repeating Characters}
\section*{\href{https://leetcode.com/problems/longest-substring-without-repeating-characters/}{题目}}
\fbox{%
  \parbox{\textwidth}{%

  \textbf{Given a string, find the length of the longest substring without repeating characters.}

  }
  }
  \section*{例子}
  \fbox{%
  \parbox{\textwidth}{%
  Input: "abcabcbb"\\
  Output: 3 \\
  Explanation: The answer is "abc", with the length of 3. \\

Input: "bbbbb"\\
Output: 1\\
Explanation: The answer is "b", with the length of 1.\\

Input: "pwwkew"\\
Output: 3\\
Explanation: The answer is "wke", with the length of 3. 
             Note that the answer must be a substring, "pwke" is a subsequence and not a substring.
  }
  }

\section*{思路}
\textbf{\color{blue}因为要找子序列，并且子序列不能有重复元素，所以一定是滑动窗口法，一个元素记录最左边+unordered\_map就可以了}
\begin{lstlisting}[caption={},frame=shadowbox]
    int lengthOfLongestSubstring(string s)
{
    if(s.empty())
        return 0;
    int left=0;
    int res=0;
    unordered_map<char,int>temp;
    for(int i=0;i<s.size();++i)
    {
        if(!temp[s[i]]||left>temp[s[i]])
        {
            res=max(res,i-left+1);
        }
        else
        {
            left=temp[s[i]];
        }
        temp[s[i]]=i+1;
    }
    return res;
}
\end{lstlisting}






\section{Add Two Numbers}
\section*{\href{https://leetcode.com/problems/add-two-numbers/}{题目} }
\fbox{%
  \parbox{\textwidth}{%

  \textbf{You are given two non-empty linked lists representing two non-negative integers. The digits are stored in reverse order and each of their nodes contain a single digit. Add the two numbers and return it as a linked list.}

  \textbf{You may assume the two numbers do not contain any leading zero, except the number 0 itself.}
  }
  }
  \section*{例子}
  \fbox{%
  \parbox{\textwidth}{%
  Input: (2 -> 4 -> 3) + (5 -> 6 -> 4)\\
  Output: 7 -> 0 -> 8\\
  Explanation: 342 + 465 = 807.
  }
  }

\section*{思路}
\textbf{\color{blue}链表的题目比较简单，一直向前遍历，直到找到结束的标志--两个链表都遍历完了以及没有进位了。}
\begin{lstlisting}[caption={},frame=shadowbox]
    ListNode* addTwoNumbers(ListNode* l1, ListNode* l2)
    {
        ListNode*temp1=l1;
        ListNode*temp2=l2;
        ListNode*start=new ListNode(0);
        ListNode*mm=start;
        int jinwei=0;
        while(temp1||temp2||jinwei)
        {
            int sum=0;
            if(temp1)
            {
                sum+=temp1->val;
                temp1=temp1->next;
            }
            if(temp2)
            {
                sum+=temp2->val;
                temp2=temp2->next;
            }
            if(jinwei)
                sum+=1;
            if(sum>=10)
                jinwei=1;
            else
                jinwei=0;
            sum=sum%10;
            ListNode*temp=new ListNode(sum);
            start->next=temp;
            start=temp;
        }
        return mm->next;
    }
\end{lstlisting}



\end{spacing}
\end{document}






